% ====================================================================
% Section 8 (v2): Discussion, limitations, future work.
% ====================================================================
\section{Discussion}
\label{sec:discussion}

\subsection{What the equivalence means}

The result of this paper is that two classical-mechanical forces---gradient
descent on a potential and a Newton's-third-law reaction to loss---reach the
operating point of purpose-built load balancers across the topologies where the
geometry leaves room to manoeuvre. We read this less as a statement about
routing and more as a statement about the problem: where the topology admits
many near-shortest alternatives, good load distribution is not a feat requiring
elaborate machinery, but something a simple local physical rule recovers on its
own. The engineering and the physics converge because they are both descending
the same landscape; the engineering simply pays for an explicit, enumerated map
of it that the physics navigates implicitly.

This reframes our own project history. The work began by asking what would
happen if a packet were a physical particle with mass, heat, and friction, and
accumulated machinery in pursuit of a decisive win. The decisive win never
materialised against strong baselines; what materialised instead, once we
stopped trying to beat the engineering and started measuring whether we could
match it, was a cleaner and more defensible claim. Parsimony was the result, not
the consolation.

\subsection{Limitations}

We state the boundaries of the claim plainly.

\begin{itemize}
\item \textbf{Simulation, not deployment.} All results are from a flow-level
simulator, not a hardware or packet-level testbed. The simulator captures
sustained-flow congestion dynamics but abstracts away queueing detail, control-
loop latency, and the operational cost of maintaining the scar state in a real
forwarding plane. In particular, micro-bursts---the microsecond-scale queue excursions that DRILL's per-packet reactivity is designed to absorb---do not exist at flow granularity; the equivalence reported against DRILL is therefore established for sustained-flow traffic and is not claimed for the micro-burst regime. A packet-level replication is the natural follow-up.
\item \textbf{Synthetic demand on real topologies.} An earlier revision of this harness loaded the G\'EANT graph under the Abilene label; the audit and the fix are reported in Section~\ref{sec:equivalence}. Both backbones now run as genuinely independent topologies, but their demand matrices remain synthetic (gravity-style) rather than measured traffic: results on them should be read as topology-faithful, demand-approximate.
\item \textbf{Predictor, not law.} The tube/sp relationship has $R^2 = 0.60$: a
strong band, not a deterministic rule, with scale-free topologies fitting
loosest.
\item \textbf{Margin choice.} The equivalence margin ($\delta = 2$ percentage
points) is a modelling decision; we report raw drop rates so the verdicts can be
reinterpreted under a different margin.
\end{itemize}

\subsection{Future work: one principle, other domains}

If load distribution on a graph is recoverable from two classical forces, the
natural question is whether the same two-term rule applies to flow problems that
are not networking at all. Last-mile logistics and supply-chain routing share
the essential structure---flows over a capacitated graph, congestion as a cost
to be spread---and admit public datasets. A companion study applying the
identical core to such a domain would test whether what we report here is a
networking trick or an instance of something more general: a classical-physics
principle for flow allocation on graphs, of which packet routing is the first
case. We regard that generalisation as the most interesting direction this work
opens, and as the proper test of whether the equivalence reported here reflects
a property of physics or merely a property of networks.

\subsection{Conclusion}

A ping packet, treated as a particle subject to a gradient and Newton's third
law, routes about as well as engineered load balancers on the topologies where
the geometry allows---and a single structural ratio says in advance which those
are. The machinery that modern load balancers bring is not, on those
topologies, buying performance the physics cannot reach; it is buying an
explicit path enumeration that the physics performs implicitly. Two terms, and a
condition for when they suffice. That is the whole of it.


Two narrower follow-ups are also queued by the analysis above. An
\emph{adaptive scar half-life}: an operator does not know failure
timescales in advance, and a slow meta-controller estimating $T_{1/2}$
from the recent variance of losses would make the temporal channel
self-tuning. And a \emph{packet-level replication} targeting the
micro-burst regime excluded above, ideally on an emulated control plane,
which would also price the telemetry this paper's simulator grants for
free.
